\documentclass[12pt]{article}
\usepackage{amsmath}
\usepackage{geometry}
\geometry{a4paper, margin=1in}

\title{Metasurface Geometry Specifications: Structural Optics for PDRC}
\author{Register 1 Blueprint}
\date{June 14, 2026}

\begin{document}

\maketitle

\section{Geometric Configuration}
The efficacy of the passive daytime radiative cooling (PDRC) metasurface relies on a specific structural geometry designed to optimize angular selectivity. Rather than relying on planar coatings, the surface utilizes a micro-scale triangular cross-section.

\subsection{Solar Reflectance via Total Internal Reflection}
The primary objective of the geometric structure is to minimize the incident solar radiation heat flux, denoted as $P_{sun}$. By engineering the apex angles of the triangular ridges, incident solar radiation ($0.3-2.5~\mu m$) undergoes total internal reflection (TIR) and is retro-reflected away from the underlying maritime infrastructure. 

\subsection{Mid-Infrared Emissivity}
Conversely, the geometry is optimized to act as a near-perfect blackbody emitter within the $8-13~\mu m$ atmospheric transparency window. The synthetic skin, composed of fluoropolymer matrices or silica-based films, leverages these structural dimensions to ensure thermal emissions are highly directional.

\subsection{Suppression of Lambertian Scattering}
By directing emissions orthogonally toward the zenith ($\theta \approx 0^\circ$), the structural geometry effectively bypasses Lambertian scattering. This prevents the re-absorption of radiated heat by the surrounding oceanic environment or adjacent vessel structures, ensuring that the net cooling power ($P_{cool}$) successfully rejects heat into deep space.

\end{document}
